\subsection{Percent Rate of Change} 
This is a piece of subtle wording that is very confusing.
\begin{itemize}
\item We said ``A value decreases by 4\% each year'' to mean that each year, the value is multiplied by \makeblank.
\item If we say ``A value decreases at a rate of 4\% per year,'' this \emph{should} mean the same thing. However, for complex reasons (which will become clear in calculus), this actually describes a ``continuous change,'' like continuous interest.
\end{itemize}
\begin{displaybox}[Continuous Growth/Decay]
Continuous growth or decay can be modeled by a function of the form
\[
f(t)=f_0 e^{rt}
\]
where $r$ is the rate of growth or decay, and $f_0$ is the value of $f(0)$. Note that this is considered \emph{decay} if the value of $f$ is decreasing; in this case, $r$ is negative.
\end{displaybox}
\begin{Example}
The cost of goods increases at a rate of 3\% per year due to inflation. If an item costs \$800 today, write a function describing its cost $c(t)$ after $t$ years.
\begin{itemize}
\item $r=\makeblanksmall=\makeblanksmall$
\item $c_0=\makeblanksmall$
\item So $c(t)=\makeblank$
\end{itemize}
\end{Example}
\begin{Example}
The amount of a drug in a patient's system decreases at a rate of 12\% per hour. If a dose of 50mg is initially administered, write a function describing the amount $d(t)$ of the drug in the patient's system after $t$ hours.
\begin{itemize}
\item $r=\makeblanksmall=\makeblanksmall$ (\makeblank[1in] because the amount of the drug is \makeblank[1in])
\item $d_0=\makeblanksmall$
\item So $d(t)=\makeblank$
\end{itemize}
\end{Example}
